\section*{Chapitre \ref{ch:iv2:commodities-and-energy} --- Matières premières et énergie}
\markboth{\protect\omTocMark{}{Réponses~: Matières premières et énergie}}{\protect\omTocMark{}{Réponses~: Matières premières et énergie}}

\begin{solution}{iq:iv2:commodities-and-energy:1}
On perd 4~\$ par baril, en situation de contango.

Le contrat à terme doit converger vers le spot à l'échéance. Si le spot reste à 70, le contrat à terme passe de 74 à 70.

\iqwrong{Dire «~le spot n'a pas bougé, donc je n'ai rien gagné~». On ne détient pas le spot. Le rendement est la variation du contrat à terme, et la pente de la courbe en fait partie.}
\end{solution}

\begin{solution}{iq:iv2:commodities-and-energy:2}
La limite haute est 88. À 92, on achète le spot, on le stocke et on vend le contrat à terme, ce qui verrouille 4~\$ par unité.

Le coût de portage complet vaut $80 \times 1.05 + 4 = 88$. On emprunte 80 et on achète~; dans un an, on livre contre le contrat à terme à 92, on rembourse 84 au titre du prêt et 4 pour le stockage.

\iqwrong{Dire que la fair value est 88. C'est un plafond. Comme l'inventaire est rare, le contrat à terme peut se situer bien en dessous.}
\end{solution}

\begin{solution}{iq:iv2:commodities-and-energy:3}
Cela vaut 25~\$ le baril. Les raffineurs sont naturellement longs.

Trois barils de brut coûtent 240. Ils donnent deux barils d'essence à $2.40 \times 42 = 100.80$ chacun, et un baril de fioul domestique à $2.70 \times 42 = 113.40$. Les produits rapportent donc 315, soit une marge de 75 sur trois barils.

\iqwrong{Oublier les unités. Les produits se cotent au gallon, le brut au baril. Il faut soustraire avant de convertir, sinon on obtient n'importe quoi.}
\end{solution}

\begin{solution}{iq:iv2:commodities-and-energy:4}
On ne gagne rien sur le mouvement parallèle, et 1\,000~\$ sur le second. On parie sur un approvisionnement plus serré, c'est-à-dire sur une backwardation plus profonde.

La position ne voit que le spread, qui passe de 2 à 3. Cela fait 1~\$ sur 1\,000 barils.

\iqwrong{Penser que cette position est longue de pétrole. Elle n'a aucune exposition directe à un mouvement parallèle, même si les contrats de première échéance ont tendance à bouger davantage que les contrats éloignés, si bien que le spread n'est pas totalement neutre en pratique.}
\end{solution}

\begin{solution}{iq:iv2:commodities-and-energy:5}
Le spark spread vaut 25~\$ par MWh, soit 40\,000~\$ pour le bloc, et le clean spark spread vaut 13.

Le combustible coûte $7 \times 5 = 35$, donc $60 - 35 = 25$. L'électricité représente $100 \times 16 = 1\,600$ MWh. Le carbone coûte $0.4 \times 30 = 12$.

\iqwrong{Y voir un profit. Cela ne tient pas compte du coût de démarrage, de la maintenance, ni du fait que la centrale peut ne pas tourner à chaque heure.}
\end{solution}

\begin{solution}{iq:iv2:commodities-and-energy:6}
Il vaut environ 11.9~\% par an.

Avec $F = S e^{(r + u - y)T}$, on a $y = r + u - \ln(F/S) = 0.06 - \ln(66/70) = 0.06 + 0.0588 = 0.1188$.

Le rendement d'opportunité (convenience yield) est la valeur de la détention du métal plutôt que d'une simple promesse de livraison~: il permet de faire tourner une usine et de honorer une commande soudaine. Il est élevé lorsque les stocks sont faibles, car une pénurie rend la dernière tonne très précieuse. Il est proche de zéro lorsque les entrepôts sont pleins.

\iqwrong{Lire la backwardation comme une prévision selon laquelle le prix va baisser de 4. Une courbe de cette forme est le prix de la détention de stock physique aujourd'hui. Elle ne dit pas, à elle seule, où ira le spot.}
\end{solution}

\begin{solution}{iq:iv2:commodities-and-energy:7}
Le contango à un mois ne peut pas dépasser environ 1.17, compte tenu de l'accès au stockage. La backwardation n'est soumise à aucune limite de ce type.

Le portage complet (full carry) vaut $50e^{0.04/12} + 1 = 51.17$. Au-delà, on achète, on stocke et on vend à terme avec un profit.

La backwardation supposerait l'inverse~: vendre le métal, acheter le contrat à terme et prendre livraison plus tard. Seul celui qui possède un inventaire peut le faire, et un détenteur qui a besoin de la matière pour faire tourner une usine ne le fera pas. Rien ne force donc le mouvement, et la courbe est plafonnée du côté du contango et libre de l'autre côté.

Il existe deux exceptions. Les limites supérieures disparaissent si les capacités de stockage sont saturées, car personne ne peut réaliser l'opération de portage. Et un produit non stockable, comme l'électricité, n'a aucun portage, si bien qu'aucune de ces bornes ne s'applique.

\iqwrong{Supposer que les deux côtés sont symétriques parce que la formule semble symétrique. Une position physique courte ne peut pas être construite à partir de rien, donc le coût de portage est une inégalité.}
\end{solution}

\begin{solution}{iq:iv2:commodities-and-energy:8}
La matière première A perd environ 11.3~\% en un an. Le gaz rapporte exactement zéro, quelle que soit la pente de sa courbe d'un mois à l'autre.

\textbf{Matière première A.} Chaque mois, le contrat que l'on détient part de 1~\% au-dessus du spot et converge vers lui. Le rendement mensuel vaut donc $1/1.01 - 1$, et l'année donne $1 - 1.01^{-12} = 11.3~\%$ de perte. Ce rendement est réel, car le spot espéré ne monte pas.

\textbf{Gaz.} Le contrat à terme de chaque mois se situe au prix de ce mois et expire à ce niveau. On ne gagne rien sur aucun contrat, puisqu'il ne bouge jamais. La pente de la courbe relève entièrement du calendrier, et acheter le contrat suivant dans une partie croissante n'est pas un coût~: il vaudra le même prix à son échéance.

Sur une année, les pentes d'une échéance à la suivante s'annulent aussi, puisque les saisons se répètent.

Ainsi, un rendement de roll naïf, calculé comme première échéance sur deuxième, qualifierait le gaz de bon portage en hiver et de mauvais en été. C'est faux~: pour mesurer le portage, il faut comparer le même mois du calendrier d'une année sur l'autre.

\iqwrong{Trier les matières premières selon la pente entre première et deuxième échéance, puis tenir pour à fort portage celles qui arrivent en tête. Sur les marchés saisonniers, on trie en réalité selon la date du calendrier.}
\end{solution}

\begin{solution}{iq:iv2:commodities-and-energy:9}
Le ratio de couverture vaut 0.686, soit environ 163 contrats, et il élimine 64~\% de la variance.

Le ratio de variance minimale vaut $h = \rho\sigma_s/\sigma_f = 0.8 \times 3/3.5 = 0.686$. Sur 10 millions de gallons, cela donne 6.857 millions de gallons de fioul domestique, soit $6\,857\,143/42\,000 = 163$ contrats.

La variance restante vaut $1 - \rho^2 = 36~\%$. L'écart-type ne baisse donc que de 40~\%. Une corrélation de 0.8 semble élevée, mais elle laisse 60~\% du risque en place.

\iqwrong{Utiliser un contrat pour 42\,000 gallons, soit un ratio de 1. On prend alors plus de risque sur le fioul domestique que la compagnie aérienne n'en a sur le kérosène, et la variance résiduelle serait supérieure à celle du minimum.}
\iqfollow{Pourquoi la compagnie aérienne couvrirait-elle son risque avec du fioul domestique~?}{Il se négocie avec une liquidité bien supérieure à celle d'un contrat sur kérosène. Le coût en est le risque de base qu'elle conserve.}
\end{solution}

\begin{solution}{iq:iv2:commodities-and-energy:10}
Le contrat à terme à un mois affiche aujourd'hui une volatilité de 50.8~\%, et celui à un an, de 8.1~\%. La moyenne pour l'option à un an est de 29.7~\%, et pour l'option à un mois de 55.3~\%.

La volatilité spot vaut $0.6$ en $\tau = 0$. On a donc $0.6\,e^{-2/12} = 0.508$ et $0.6\,e^{-2} = 0.081$.

Une option sur contrat à terme à l'échéance $T$ fait intervenir la variance moyenne lorsque $\tau$ décroît de $T$ à 0~:
\[
\bar\sigma^2 = \frac{0.36}{T}\int_0^T e^{-4\tau}\,\mathrm{d}\tau = 0.36\,\frac{1 - e^{-4T}}{4T}.
\]
Pour $T = 1$, $\bar\sigma = 0.6 \times 0.495 = 0.297$. Pour $T = 1/12$, $\bar\sigma = 0.553$.

C'est l'\index{effet Samuelson}effet Samuelson. Un choc sur le spot est largement résorbé dans le temps par les réponses de l'offre, si bien qu'un contrat éloigné réagit peu. À mesure qu'il approche de l'échéance, il ressent chaque choc. Cela rend visibles, dans la courbe des volatilités implicites, le taux de retour à la moyenne et le niveau de volatilité.

\iqwrong{Retenir la volatilité actuelle du contrat à terme à un an, soit 8.1~\%, pour l'option à un an. Au fil de l'année, ce contrat vieillit et devient celui de première échéance, si bien que ce chiffre est bien trop faible.}
\end{solution}

\begin{solution}{iq:iv2:commodities-and-energy:11}
L'option moyennée vaut environ 1.86, contre environ 3.22 pour l'option européenne. Elle coûte donc environ 58~\% de l'option européenne.

Moyenner sur toute la durée de vie divise la variance effective par trois. La volatilité effective vaut donc $40~\%/\sqrt{3} = 23.1~\%$ à même échéance. Avec l'approximation à la monnaie, $0.4 \times 70 \times 0.231 \times \sqrt{1/12} = 1.87$, et la formule de Black donne 1.86. L'option classique vaut $0.4 \times 70 \times 0.40 \times \sqrt{1/12} = 3.23$, et Black donne 3.22.

Les consommateurs de matières premières achètent et vendent le physique tout au long du mois. Un prix moyen est donc ce qu'ils paient réellement. Et il rend le règlement-livraison difficile à manipuler, car personne ne peut faire basculer le résultat en agissant sur un seul relevé.

\iqwrong{Utiliser $\sigma/2$ ou $\sigma/\sqrt{2}$. La variance d'une moyenne cumulée vaut $\sigma^2T/3$, ce qui découle de l'intégration de $(T - t)^2/T^2$.}
\iqfollow{Et si la moyenne portait sur 22 fixings quotidiens~?}{Le facteur de variance vaut $(n + 1)(2n + 1)/(6n^2) = 0.356$, légèrement supérieur à $1/3$, si bien que la volatilité est d'environ 23.9~\%.}
\end{solution}

\begin{solution}{iq:iv2:commodities-and-energy:12}
L'écart est tout juste ouvert~: environ 20 cents le baril, avant la qualité, le port et les autres frais.

Vente moins achat moins fret~: $77 - 73 - 3.50 = 0.50$. Financer la cargaison coûte $73 \times 0.05 \times 30/365 = 0.30$, ce qui laisse 0.20.

Pour la figer, on achète la cargaison et on vend le benchmark européen pour la date de livraison. Le risque de prix américain disparaît dès l'achat. Le fret se fixe en réservant le navire ou par un contrat de fret.

\textbf{Ce qui reste.}
\begin{itemize}
\item Qualité~: la nuance de brut ne se négocie peut-être pas au prix du benchmark.
\item L'écart du benchmark par rapport à la vente effective peut évoluer.
\item Calendrier~: retard, surestaries et météo. La date de la couverture est fixée, mais l'arrivée du navire ne l'est pas.
\item Exécution~: si de nombreux navires se dirigent vers le même port, le prix local à la livraison baisse. L'arbitrage se referme de lui-même.
\end{itemize}

\iqwrong{Prendre pour un profit un spread entre deux contrats à terme. Les deux contrats portent sur des lieux différents. L'écart ne signifie rien tant qu'on n'a pas payé pour déplacer le baril.}
\end{solution}

\begin{solution}{iq:iv2:commodities-and-energy:13}
Les prix valent 0~\$, 10~\$, 40~\$ et 90~\$ pour 20, 60, 100 et 115 MW. À 140, la demande ne peut pas être couverte, si bien que le prix va jusqu'au plafond applicable, ou qu'une partie de la charge est coupée. Avec 60 MW d'éolien, les prix tombent à 0~\$, 0~\$, 10~\$ et 40~\$.

Ranger les centrales par coût marginal, et la centrale marginale fixe le prix. La capacité cumulée jusqu'au point de demande~: 30, 70, 110, 130. Ainsi, 20 MW se situent dans l'éolien, 60 dans le nucléaire, 100 dans le gaz et 115 dans les centrales de pointe.

Avec 30 MW supplémentaires à coût nul, tout l'empilement se décale vers la droite~: 60, 100, 140, 160. Ainsi, 100 tombe dans le nucléaire et 115 dans le gaz. C'est l'\index{effet d'ordre de préséance économique}effet d'ordre de préséance économique (merit-order effect). Une offre à coût nul supplémentaire fait baisser les prix en général, y compris ceux que perçoit l'éolien lui-même.

\textbf{Prix négatifs.} L'électricité ne se stocke pas à grande échelle, et l'offre doit égaler la demande à chaque instant. Certaines centrales paient pour rester en service, car l'arrêt puis le redémarrage coûtent cher. Et les éoliennes et centrales solaires subventionnées peuvent gagner davantage que le prix négatif, si bien qu'elles continuent de produire. Elles offrent donc en dessous de zéro pour rester en service, et le prix se fixe à ce niveau.

\begin{dated}{2026-11}{Heures négatives en Allemagne}\label{dat:iv2:commodities-and-energy:neg}
Le régulateur allemand a compté des prix de gros négatifs pendant 573 heures en 2025 et 457 heures en 2024. Il faut consulter le dernier décompte avant d'en citer un.
\end{dated}

\iqwrong{Faire la moyenne des coûts des centrales. Sur un marché à prix marginal, chaque producteur reçoit le prix de l'unité marginale, et c'est ainsi que les centrales éoliennes et nucléaires couvrent leurs coûts fixes.}
\iqfollow{Un forward en base vaut 50 et celui de pointe 70, avec 80 heures de pointe et 88 heures creuses par semaine. Quel est le prix hors pointe~?}{$(168 \times 50 - 80 \times 70)/88 = 31.82$.}
\end{solution}

\begin{solution}{iq:iv2:commodities-and-energy:14}
L'option vaut environ 6.94 par MWh, dont 5 de valeur intrinsèque et 1.94 de valeur temps. Pour une corrélation de 0.9, elle tombe à environ 5.56. Pour la couvrir, on vend environ 0.74 MWh d'électricité et on achète environ 4.52 MMBtu de gaz.

Soit la jambe du strike $F_2 = 7 \times 5 = 35$. Alors la formule de \index{formule de Margrabe}Margrabe utilise la volatilité du ratio,
\[
\sigma = \sqrt{0.6^2 + 0.4^2 - 2 \times 0.5 \times 0.6 \times 0.4} = 0.529,
\]
avec $d_1 = (\ln(40/35) + \sigma^2 T/2)/(\sigma\sqrt{T}) = 0.637$ et $d_2 = 0.372$. Donc $V = 40\,\Phi(d_1) - 35\,\Phi(d_2) = 6.94$.

Pour une corrélation de 0.9, la volatilité du ratio vaut 0.297, si bien que la valeur est surtout intrinsèque. Et quand la corrélation tend vers un, à volatilités égales, elle est exactement intrinsèque.

\textbf{Couverture.} Les deltas valent $\Phi(d_1) = 0.738$ sur l'électricité et $\Phi(d_2) = 0.645$ sur la jambe gaz. Cela fait $0.645 \times 7 = 4.52$ MMBtu de gaz par MWh. Comme pour toute option dont on est long, on réajuste la couverture en vendant haut et en achetant bas, ce qui rapporte la volatilité réalisée du spread.

\iqwrong{Se fonder sur la seule volatilité de l'électricité. Ce qui compte, c'est la volatilité du spread, si bien qu'une forte corrélation rend la centrale moins chère.}
\iqfollow{Que vaut-elle si l'électricité tombe à 30~?}{Environ 1.48~: une centrale hors de la monnaie aujourd'hui vaut encore quelque chose, car elle peut attendre, et c'est pourquoi on ne la ferme pas au premier mauvais mois.}
\end{solution}

\begin{solution}{iq:iv2:commodities-and-energy:15}
Rien ne l'en empêche. Quand les capacités de stockage sont saturées, la borne de portage ne tient plus, et un détenteur qui ne peut pas prendre livraison doit payer quelqu'un pour reprendre le contrat. Un prix de $-38$ alors que le mois suivant vaut 20 signifie que le stockage pour le mois vaut environ 58~\$ le baril.

Si l'on peut stocker, $F_1 \ge F_2 - c$ est vérifiée~: à un prix inférieur, on achète, on stocke un mois et on vend le contrat suivant. La borne inférieure du contrat le plus proche est donc le contrat suivant moins le coût de portage. Pour un coût de stockage $c$ de 58~\$, $20 - 58 = -38$.

Normalement, $c$ est un petit nombre. Quand la capacité est pleine, ce n'est plus un prix que l'on peut payer, si bien que la borne n'existe pas. Le contrat le plus proche est donc libre de se fixer là où le marché le place.

Trois catégories de détenteurs sont contraintes de vendre. Un fonds long only, y compris les produits de détail, qui doit sortir à sa date de roll. Un spéculateur qui n'a aucun moyen de prendre livraison. Et toute entité dont les limites de risque imposent de clôturer la position.

Les acheteurs capables d'assumer la livraison sont peu nombreux et détiennent des capacités de stockage. Ce sont eux qui fixent le prix.

\begin{dated}{2026-11}{Le jour où c'est arrivé}\label{dat:iv2:commodities-and-energy:neg-wti}
Le 20 avril 2020, le contrat WTI de mai, arrivant à échéance, a fixé son prix de règlement à $-\$37.63$, le contrat suivant s'établissant à 20.43, soit un calendar spread d'environ $-\$58$. Le deuxième écart le plus large entre les deux premières échéances dans la série quotidienne publique de 1985 à 2024 est de 8.49~\$.
\end{dated}

\iqwrong{Dire qu'un prix ne peut pas passer sous zéro parce que le pétrole a de la valeur. Le contrat à terme est un droit à prendre livraison en un lieu et à une date donnés. Quand on ne peut pas en prendre livraison, l'obligation a une valeur négative.}
\iqfollow{Pourquoi les autres contrats n'ont-ils pas été touchés aussi durement~?}{Ils avaient du temps. Le deuxième mois disposait d'un mois entier pour écouler le stock, et les producteurs pouvaient réduire leur production.}
\end{solution}

\begin{solution}{iq:iv2:commodities-and-energy:16}
La valeur intrinsèque vaut 0.72~: injecter au mois 3 et soutirer au mois 5. Avec l'information révélée au mois 4, le réservoir vaut 0.87, si bien que l'option d'attendre pour agir ajoute 0.15.

\textbf{Intrinsèque.} On cherche la meilleure paire de dates, on achète à la date la plus tôt et on vend à la plus tardive. La paire retenue est le mois 3 à 2.85 et le mois 5 à 3.60, les coûts totalisent 0.03, soit $3.60 - 2.85 - 0.03 = 0.72$. Un seul cycle est possible~: le réservoir ne contient qu'une unité, et aucune autre paire n'ajoute de valeur. L'énumération exhaustive le confirme.

\textbf{Extrinsèque.} On injecte au mois 3. Au mois 4, on observe le prix du mois 5, et l'énumération de chaque action dans chaque état confirme le plan ci-dessous.
\begin{itemize}
\item À 4.10, on attend~: on obtient $4.10 - 2.85 - 0.03 = 1.22$.
\item À 3.10, on soutire au prix du mois 4, soit 3.40~: on obtient $3.40 - 2.85 - 0.03 = 0.52$.
\end{itemize}
Donc la valeur vaut $(1.22 + 0.52)/2 = 0.87$. Cela revient à un plancher à 3.40 sur la vente, et le forward à 3.60 reste inchangé en espérance.

En pratique, on verrouille la valeur intrinsèque avec une couverture statique, puis on fait rouler la couverture à mesure que la courbe bouge, en basculant chaque fois vers un meilleur spread. Chaque bascule verrouille davantage de valeur, et c'est ce que l'on appelle l'\index{intrinsèque glissant}intrinsèque glissant (rolling intrinsic). Elle ne compte pas toute la valeur, mais elle donne une borne inférieure qu'on peut sécuriser.

\iqwrong{Ne présenter que la valeur intrinsèque. Un réservoir est un ensemble d'options sur spread portant sur la courbe à terme, et une grande partie de la valeur d'un réservoir réel peut tenir à sa flexibilité. Si rien ne bouge, la valeur extrinsèque est nulle, mais si la volatilité est forte, elle peut en constituer la part principale.}
\iqfollow{Qu'est-ce qui limite un réservoir réel~?}{Le rythme d'injection dépend du niveau de remplissage, et le gaz exige un niveau minimal de stock utile. Ces limites transforment le problème en un programme dynamique dont le remplissage est une variable d'état.}
\end{solution}

\begin{solution}{iq:iv2:commodities-and-energy:17}
L'appel de marge s'élève à 200~millions de dollars, et la valeur actuelle du gain sur les ventes vaut environ 157~millions de dollars. La couverture un pour un est donc trop grande, même dans le meilleur cas. Une couverture ajustée (tail hedge) serait d'environ 79~\%.

\textbf{Mouvement plat.} Les barils de chaque année rapportent 20~millions de dollars, somme qu'on ne perçoit qu'à la livraison. À 5~\%, les flux arrivant en milieu d'année, le total est~:
\[
\sum_{k = 0}^{9} 20 \times e^{-0.05(k + 0.5)} = \$157.4 \text{ millions de dollars},
\]
contre 200~millions de dollars payés aujourd'hui. On est donc en perte de 42.6~millions de dollars en valorisation au prix de marché, et l'appel de marge tombe immédiatement alors que les gains s'étalent sur une décennie.

\textbf{Ajuster la couverture.} Mettre la pile à l'échelle de $157.4/200 = 78.7~\%$ la rend neutre à ce mouvement. Le problème de trésorerie est alors plus petit, mais il subsiste~: on paie encore 157~millions de dollars aujourd'hui contre des gains à venir. L'entreprise doit disposer de cette ligne de crédit, signée à l'avance.

\textbf{L'extrémité longue de la courbe bouge de 30~\%.} Les prix à longue échéance bougent en général moins que ceux des échéances proches. Si la vente moyenne ne rapporte que 6~\$ au lieu de 20~\$, on gagne environ $0.3 \times 157 = \$47$ millions sur les ventes, et l'on perd toujours 200~millions de dollars sur la pile. Un ratio de couverture issu d'un mouvement plat est alors largement trop grand~: le bon ratio doit provenir de la façon dont chaque partie de la courbe bouge réellement.

\textbf{Contango.} À 50 cents par mois sur 10~millions de barils, on perd 5~millions de dollars à chaque roll, soit environ 60~millions de dollars par an au départ, somme qui diminue à mesure que la pile se réduit.

Une même position peut donc avoir une valeur économique sûre et mourir pourtant faute de financement. Un prêteur qui voit des appels de marge la première année ne voit pas les gains de la dixième.

\iqwrong{Dire «~Être long de pétrole contre une vente à découvert de pétrole, donc c'est couvert.~» Le risque vient de trois choses~: la courbe peut changer de forme, les flux de trésorerie arrivent à des dates différentes, et le coût du roll dépend de la forme de la courbe.}
\end{solution}

\begin{solution}{iq:iv2:commodities-and-energy:18}
À volatilité plate, le prix d'exercice du call est d'environ 84.30. Avec le skew, il tombe à environ 75.43.

\textbf{Plat.} Sous Black avec des taux nuls, le put à 60 et 30~\% coûte 3.74. On résout le prix d'exercice du call avec la même prime~: 84.30. Le revenu du producteur est donc garanti au minimum à 60 et plafonné à 84.30.

\textbf{Skew.} Le put à 36~\% coûte 5.11, et un call à 26~\% doit encaisser 5.11, d'où un prix d'exercice de 75.43. La structure protège le même plancher, mais abandonne près de neuf dollars de hausse, car le producteur vend de la volatilité bon marché et en achète à prix élevé.

Le producteur renonce à la hausse au-delà du prix d'exercice du call. «~Sans coût~» signifie qu'aucune somme n'est versée d'avance. Le coût se trouve dans cette hausse abandonnée, auquel s'ajoute la marge du dealer intégrée aux prix d'exercice.

À la limite, un put et un call de même prix d'exercice $K = F$ donnent un forward synthétique vendeur, à prime nulle.

\iqwrong{Assimiler une prime nulle à un coût nul. Tout collar au juste prix a une espérance de gain nulle avant la marge du dealer. Et lorsque le skew joue contre vous, les conditions se dégradent sans qu'aucune prime n'apparaisse dans la transaction.}
\iqfollow{Que fait la banque de son côté~?}{Elle détient le put vendu et le call acheté, elle est donc longue du delta du contrat à terme. Elle vend des contrats à terme pour couvrir ce delta, et gère l'exposition au skew qu'elle a reprise de nombreuses opérations de ce type.}
\end{solution}

\begin{solution}{iq:iv2:commodities-and-energy:19}
L'approximation de Kirk donne environ 7.46, et un modèle de spread normal environ 7.70. Une simulation donne 7.47. C'est donc l'approximation de Kirk qu'il faut retenir ici, la simulation servant à la vérifier.

\textbf{Kirk.} On traite $F_2 + K$ comme un actif lognormal unique, dont la volatilité est pondérée par $w = F_2/(F_2 + K) = 80/95 = 0.842$. Alors
\[
\sigma = \sqrt{0.35^2 + (0.842 \times 0.3)^2 - 2 \times 0.9 \times 0.35 \times 0.842 \times 0.3} = 0.165.
\]
Et $V = F_1\Phi(d_1) - (F_2 + K)\Phi(d_2)$, avec $d_1, d_2$ issus de $\ln(100/95)$. On obtient 7.46. Pour $K = 0$, la formule est exacte~: c'est la formule de Margrabe.

\textbf{Spread normal.} Les volatilités en dollars valent $35$ et $24$, celle du spread vaut donc~:
\[
\sqrt{35^2 + 24^2 - 2 \times 0.9 \times 35 \times 24} = 17
\]
par an, et $12.02$ sur six mois. Avec un spread à terme $x = 5$,
\[
V = 12.02\,\bigl(\varphi(5/12.02) + (5/12.02)\Phi(5/12.02)\bigr) = 7.70.
\]
Le modèle s'écarte d'environ 3~\% ici, car le vrai spread entre deux lognormales n'est pas normal.

\textbf{Contrôle.} Simuler les deux contrats à terme avec la corrélation donnée, et prendre la moyenne du gain. L'approximation de Kirk perd en précision lorsque les prix d'exercice sont élevés par rapport à $F_2$, ou lorsque la corrélation est très forte.

\iqwrong{Utiliser une seule volatilité de spread, issue du premier modèle, pour tous les prix d'exercice. La corrélation qui reproduit les options sur spread du marché varie avec le prix d'exercice~: le skew de corrélation est donc un risque réel dans ces livres, et les spreads à faible volatilité le masquent.}
\end{solution}

\begin{solution}{iq:iv2:commodities-and-energy:20}
La valeur intrinsèque est de 70\,400~\$. Comme décision mensuelle, elle vaut environ 108\,000~\$, et comme 22 décisions quotidiennes, environ 196\,000~\$. Je paierais nettement moins que le chiffre quotidien, au plus environ 130\,000 à 160\,000~\$, pour les raisons exposées ci-dessous.

\textbf{Intrinsèque.} La jambe à prix d'exercice vaut $8 \times 5 = 40$, le spread est donc de 2~\$. Il y a $100 \times 16 \times 22 = 35\,200$~MWh. Donc $2 \times 35\,200 = \$70\,400$.

\textbf{Mensuel.} Margrabe sur 42 contre 40, avec une volatilité du ratio
\[
\sqrt{0.25 + 0.1225 - 2 \times 0.6 \times 0.5 \times 0.35} = 0.403
\]
Avec $T = 1/12$, on obtient 3.07~\$ par MWh, soit 108\,000~\$.

\textbf{Quotidien.} Chaque jour, la volatilité du ratio vaut $\sqrt{0.81 + 0.25 - 2 \times 0.5 \times 0.9 \times 0.5} = 0.781$, pour $T = 1/8$. On obtient 5.57~\$ par MWh, soit 196\,000~\$.

La valeur quotidienne est plus élevée pour deux raisons. On peut sauter les jours où le spread est négatif, si bien qu'on détient plusieurs options et non une seule. Et les prix spot sont plus volatils que les contrats à terme mensuels.

\textbf{Couvertures.} Face à la valeur mensuelle, il faut vendre $0.683 \times 35\,200 = 24\,059$~MWh d'électricité et acheter $0.641 \times 8 \times 35\,200 = 180\,501$~MMBtu de gaz. Face à la valeur quotidienne, les deltas sont plus faibles~: 21\,948~MWh et 145\,138~MMBtu. Ces deux montants restent inférieurs à une position pleine, car la centrale ne tourne pas dans tous les états.

\textbf{Les frais.} Le chiffre quotidien suppose qu'on puisse faire tourner ou arrêter la centrale sans frais, et capter le prix de chaque jour. Face à cela~:
\begin{itemize}
\item Chaque démarrage coûte du combustible et de l'usure, et les durées minimales de fonctionnement empêchent de ne retenir que les bonnes heures.
\item La volatilité spot n'est pas monétisée par une couverture en delta au moyen de contrats à terme mensuels. On détient l'écart entre exposition mensuelle et exposition quotidienne comme un risque, et le couvrir exige des produits quotidiens liquides.
\item Des pannes peuvent laisser la position vendeuse en électricité le pire jour, précisément quand la couverture a le plus besoin de la centrale.
\item L'estimation horaire ou quotidienne de la corrélation est très incertaine, et passer de 0.5 à 0.6 modifie fortement la valeur.
\end{itemize}
Je ne paierais donc au-dessus du chiffre mensuel que ce que je peux couvrir ou ce qui me paraît fondé, et je laisserais une marge nette sous le haut de la fourchette. Les chiffres présentés ici illustrent une manière de raisonner, et non des règles.

\textbf{Forward à 32~\$.} Le spread vaut alors $-8$, la valeur intrinsèque est nulle. La valeur mensuelle s'effondre à environ 0.04~\$ par MWh, soit environ 1\,500~\$. La valeur quotidienne vaut environ 1.16~\$, soit environ 41\,000~\$. Très hors de la monnaie, presque toute la valeur tient à la volatilité d'un jour à l'autre. Une centrale qui paraît morte sur la courbe mensuelle conserve donc une vraie valeur grâce aux pics de prix.

\iqwrong{Payer la valeur d'option calculée à partir d'une seule volatilité à terme pour les 22 jours. Cela ne tient compte ni de la différence entre décisions mensuelles et quotidiennes, ni des limites d'exploitation.}
\iqfollow{Qui est votre contrepartie naturelle sur cette opération~?}{Un propriétaire de centrale qui veut un revenu fixe et aucun risque de marché. On achète son option, et il vend sa volatilité contre du cash.}
\end{solution}
